\documentclass{article}
\usepackage{amsmath, amssymb, amsthm}
\usepackage{hyperref}
\usepackage{geometry}
\geometry{margin=1in}

\title{Erd\H{o}s Problem \#382}
\date{Last updated: 2025-08-31}
\author{Source: erdosproblems.com}

\begin{document}
\maketitle

\noindent\textbf{Status:} OPEN \quad \textbf{No prize}

\noindent\textbf{Tags:} number theory

\medskip

\noindent\textbf{Problem Statement:}

Let $u\leq v$ be such that the largest prime dividing $\prod_{u\leq m\leq v}m$ appears with exponent at least $2$. Is it true that $v-u=v^{o(1)}$? Can $v-u$ be arbitrarily large?




Erd\H{o}s and Graham report it follows from results of Ramachandra that $v-u\leq v^{1/2+o(1)}$. Cambie has observed that the first question boils down to some old conjectures on prime gaps. By Cram\'{e r's conjecture}, for every $\epsilon&gt;0,$ for every $u$ sufficiently large there is a prime between $u$ and $u+u^\epsilon$. Thus for $u+u^\epsilon&lt;v$, the largest prime divisor of \( \prod_{u \leq m \leq v} m \) appears with exponent $1$. Since this is not the case in the question, \( v - u = v^{o(1)} \). Cambie also gives the following heuristic for the second question. The 'probability' that the largest prime divisor of $n$ is $&lt;n^{1/2}$ is $1-\log 2&gt;0$. For any fixed $k$, there is therefore a positive 'probability' that there are $k$ consecutive integers around $q^2$ (for a prime $q$) all of whose prime divisors are bounded above by $q$, when $v-u\geq k$. See [383] for a conjecture along these lines. A similar argument applies if we replace multiplicity $2$ with multiplicity $r$, for any fixed $r\geq 2$. See also [380] .

\noindent\textbf{Related OEIS sequences:} A388850


\bigskip
\noindent\small{Source: \url{https://www.erdosproblems.com/382}}
\end{document}
